% =========================================================================
% DOCUMENTCLASS
% =========================================================================
\documentclass[12pt,oneside]{book}

% =========================================================================
% METADATOS
% =========================================================================
\newcommand{\universidad}{Universidad de Buenos Aires (UBA)}
\newcommand{\facultad}{Facultad de Derecho}
\newcommand{\carrera}{Carrera de Abogacía}
\newcommand{\materia}{Derecho de la Integración}
\newcommand{\claseinfo}{Clase de repaso integrador --- Ciclo lectivo 2026}
\newcommand{\titulo}{Jerarquía normativa, derecho comunitario, fuentes derivadas y armonización legislativa en los procesos de integración}
\newcommand{\autor}{Isaac Edgar Camacho Ocampo}
\newcommand{\anio}{2026}

% =========================================================================
% PAQUETES
% =========================================================================
\usepackage[utf8]{inputenc}
\usepackage[T1]{fontenc}
\usepackage[provide=*,spanish]{babel}

\usepackage[
    a4paper,
    top=2.5cm,
    bottom=2.5cm,
    left=3cm,
    right=2.5cm,
    headheight=15pt
]{geometry}

\usepackage{amsmath}
\usepackage{amssymb}
\usepackage{mathtools}

\usepackage{graphicx}
\usepackage{float}
\usepackage{caption}
\usepackage{subcaption}

\usepackage{booktabs}
\usepackage{array}
\usepackage{longtable}
\usepackage{tabularx}

\usepackage{enumitem}

\usepackage{xcolor}
\usepackage[most]{tcolorbox}

\usepackage{fancyhdr}
\usepackage{ragged2e}

\graphicspath{
    {./img/}
    {./graficos/}
    {./figuras/}
}

% =========================================================================
% CONFIGURACIÓN
% =========================================================================
\setcounter{tocdepth}{2}
\setcounter{secnumdepth}{0}

\setlength{\parindent}{0pt}
\setlength{\parskip}{0.5em}

\setlist[itemize]{topsep=0.3em, itemsep=0.2em}
\setlist[enumerate]{topsep=0.3em, itemsep=0.2em}

\clubpenalty=10000
\widowpenalty=10000

\addto\captionsspanish{\renewcommand{\chaptername}{Bloque}}

\pagestyle{fancy}
\fancyhf{}
\fancyhead[L]{\nouppercase{\leftmark}}
\fancyhead[R]{\materia}
\fancyfoot[C]{\thepage}
\renewcommand{\headrulewidth}{0.4pt}
\renewcommand{\footrulewidth}{0pt}

\fancypagestyle{plain}{
    \fancyhf{}
    \fancyfoot[C]{\thepage}
    \renewcommand{\headrulewidth}{0pt}
}

% =========================================================================
% COMANDOS
% =========================================================================
\newcommand{\abs}[1]{\left|#1\right|}
\newcommand{\norm}[1]{\left\lVert#1\right\rVert}
\newcommand{\vect}[1]{\mathbf{#1}}
\newcommand{\deriv}[2]{\frac{d#1}{d#2}}
\newcommand{\pderiv}[2]{\frac{\partial#1}{\partial#2}}

% =========================================================================
% ENTORNOS / CAJAS
% =========================================================================
\newtcolorbox{definition}[1]{
    enhanced, breakable,
    title={Definición: #1},
    fonttitle=\bfseries,
    colback=blue!3, colframe=blue!50!black,
    boxrule=0.8pt, arc=2mm
}

\newtcolorbox{important}{
    enhanced, breakable,
    title={Importante},
    fonttitle=\bfseries,
    colback=yellow!8, colframe=orange!70!black,
    boxrule=0.8pt, arc=2mm
}

\newtcolorbox{example}[1]{
    enhanced, breakable,
    title={Ejemplo: #1},
    fonttitle=\bfseries,
    colback=green!4, colframe=green!40!black,
    boxrule=0.8pt, arc=2mm
}

\newtcolorbox{formula}[1]{
    enhanced, breakable,
    title={Fórmula / Regla: #1},
    fonttitle=\bfseries,
    colback=purple!4, colframe=purple!50!black,
    boxrule=0.8pt, arc=2mm
}

\newtcolorbox{exercise}[1]{
    enhanced, breakable,
    title={Ejercicio: #1},
    fonttitle=\bfseries,
    colback=gray!5, colframe=gray!60!black,
    boxrule=0.8pt, arc=2mm
}

\newtcolorbox{note}{
    enhanced, breakable,
    title={Nota},
    fonttitle=\bfseries,
    colback=white, colframe=gray!50,
    boxrule=0.6pt, arc=2mm
}

% =========================================================================
% CONFIGURACIÓN FINAL
% =========================================================================
\usepackage[
    colorlinks=true,
    linkcolor=blue!50!black,
    citecolor=green!40!black,
    urlcolor=blue!60!black,
    pdftitle={\titulo},
    pdfauthor={\autor},
    pdfsubject={Apuntes universitarios},
    bookmarks=true
]{hyperref}

% =========================================================================
% DOCUMENT
% =========================================================================
\begin{document}

% ---------------------------------------------------------------
% PORTADA
% ---------------------------------------------------------------
\begin{titlepage}
\thispagestyle{empty}

\begin{center}

\vspace*{1cm}

{\large \textsc{\universidad}}

\vspace{0.2cm}

{\large \textsc{\facultad}}

\vspace{0.2cm}

{\large \carrera}

\vspace{0.3cm}

{\normalsize \claseinfo}

\vspace{2.5cm}

{\Huge \textbf{\materia}}

\vspace{1cm}

{\LARGE \titulo}

\vspace{0.8cm}

\textit{Comprender cada instituto es el primer paso hacia el criterio jurídico propio.}

\vspace{1.2cm}

\rule{0.70\textwidth}{0.5pt}

\vspace{0.5cm}

{\large Apuntes de estudio}

\vspace{0.5cm}

\rule{0.70\textwidth}{0.5pt}

\vfill

{\large \autor}

\vspace{0.3cm}

{\large \anio}

\end{center}

\vfill

\begin{flushleft}
\textit{Este es un modesto aporte a los estudiantes de ``Derecho de la Integración'' no pretende sustituir el material oficial de la materia.}
\end{flushleft}

\end{titlepage}

% ---------------------------------------------------------------
% ÍNDICE
% ---------------------------------------------------------------
\tableofcontents

% =========================================================================
% CONTENIDO
% =========================================================================

\chapter*{Introducción}
\addcontentsline{toc}{chapter}{Introducción}
\markboth{Introducción}{Introducción}

El presente documento reconstruye una clase de repaso integrador de Derecho de la Integración, organizada en tres grandes bloques temáticos que se corresponden con los capítulos de este material: los fundamentos del derecho comunitario, las fuentes derivadas y los procedimientos de armonización legislativa, y las instituciones y tensiones actuales del proceso de integración, con especial referencia al Mercosur.

\section*{Temas tratados en esta clase}
\begin{itemize}
    \item Jerarquía normativa entre tratados y leyes en la Argentina, y el alcance del artículo 75, inciso 24, de la Constitución Nacional.
    \item Características del derecho comunitario: especificidad, autonomía, supremacía, aplicabilidad inmediata y efecto directo (vertical y horizontal).
    \item Tipos de normas derivadas de la Unión Europea: reglamento, decisión y directiva.
    \item Procedimientos de armonización legislativa: uniformidad, armonía y aproximación, con el caso aplicado del Reglamento europeo de insolvencia.
    \item Composición y sistema de votación del Parlamento Europeo, en comparación con el Parlamento del Mercosur.
    \item Tensión entre integración y soberanía, y clasificación del Mercosur como unión aduanera imperfecta o incompleta.
\end{itemize}

\section*{Utilidad profesional de estos contenidos}

Estos contenidos resultan relevantes para el ejercicio profesional en comercio exterior, contratos internacionales, aduanas, arbitraje, derecho migratorio o asesoramiento a empresas que operan en el Mercosur, en la medida en que permiten determinar qué norma prevalece cuando un tratado de integración entra en conflicto con una ley interna, y qué efectos concretos produce esa prevalencia sobre un contrato, una importación o una situación migratoria.

Distinguir entre reglamento, decisión y directiva, así como entre uniformidad y armonía legislativa, permite anticipar cómo se resolverá un conflicto normativo antes de litigarlo: por ejemplo, si un cliente exportador puede invocar directamente un derecho frente a un Estado, o si primero debe reclamar que ese Estado incorpore una norma que todavía no aplicó.

Más allá del ejercicio profesional, comprender estos mecanismos permite abordar con criterio propio los debates actuales sobre soberanía, migración y globalización, cuestiones que exceden el ámbito académico y atraviesan la vida ciudadana en un mundo cada vez más interconectado.

\begin{example}{La analogía del edificio compartido}
El proceso de integración puede compararse con un edificio de propietarios. El Bloque 1 corresponde a los cimientos: las reglas constitucionales que definen qué jerarquía tiene el reglamento del consorcio (el derecho comunitario) frente a las reglas de cada departamento (el derecho interno de cada Estado). El Bloque 2 representa el reglamento interno del edificio en funcionamiento: quién dicta qué tipo de norma (una decisión de la asamblea, una circular general, una directiva que cada propietario aplica a su manera) y cómo se concilian costumbres distintas. El Bloque 3 es la asamblea de propietarios en acción: cómo se reparten los votos, qué tensiones surgen cuando alguien siente que cedió demasiada autonomía, y qué tan terminado --- o inconcluso --- está en la práctica ese edificio (el caso del Mercosur).
\end{example}

% =========================================================================
% BLOQUE 1
% =========================================================================
\chapter{Fundamentos del derecho comunitario}

\section{Unidad 1 --- Jerarquía normativa entre tratados y leyes en la República Argentina}

En la República Argentina, la jerarquía normativa entre los tratados y las leyes internas atravesó una evolución concreta. A partir del caso jurisprudencial conocido como ``Cafés La Virginia'', la Corte Suprema de Justicia de la Nación reconoció a los tratados de integración la misma jerarquía normativa que ostentan los demás tratados internacionales frente a la ley. Posteriormente, la reforma constitucional de 1994 consolidó esa prelación normativa al incorporarla expresamente en el artículo 75, inciso 24, de la Constitución Nacional.

\begin{formula}{Constitución Nacional Argentina, art. 75, inciso 24}
``Corresponde al Congreso: [\ldots] Aprobar tratados de integración que deleguen competencias y jurisdicción a organizaciones supraestatales en condiciones de reciprocidad e igualdad, y que respeten el orden democrático y los derechos humanos. Las normas dictadas en su consecuencia tienen jerarquía superior a las leyes.''
\end{formula}

Corresponde señalar una advertencia central sobre este punto, ya que no debe confundirse la solución del derecho argentino con una regla universal del derecho de la integración.

\begin{important}
\textit{``Esa es la postura de la República Argentina. No significa que el derecho comunitario per se tenga jerarquía frente al derecho interno.''}

Es decir: que un Estado reconozca jerarquía superior a las normas de un proceso de integración es una decisión de su propio derecho interno (en la Argentina, de rango constitucional), y no una característica automática y universal de todo derecho de la integración frente a cualquier ordenamiento nacional.
\end{important}

\section{Unidad 2 --- El doble control: procedimiento de vigencia simultánea}

La normativa de ``fuente derivada'' --- es decir, las normas dictadas por los órganos de un proceso de integración, más allá de que deban expedirse por consenso --- entra en vigor mediante una doble instancia de control, denominada ``doble chequeo'':

\begin{itemize}
    \item Una vez emitida la norma por el órgano competente, debe ser notificada --- a través de la secretaría del organismo --- a todos los Estados parte.
    \item Cada Estado, según lo disponga su propio sistema jurídico interno, debe proceder a la incorporación de esa norma derivada a su derecho interno, y comunicar ese acto de incorporación a la secretaría.
    \item Recién cuando la secretaría recibe la totalidad de los documentos de incorporación de todos los Estados, informa a todos ellos la fecha en que la norma entra en vigencia.
\end{itemize}

\begin{quote}
\itshape ``Hay un doble chequeo por parte del Estado, porque no solamente se necesita la conformidad para que la norma se dicte, sino que hasta que vos no la incorporás y no lo informás, la norma no puede entrar en vigor para nadie.''
\end{quote}

Este mecanismo de doble control (propio, entre otros, del sistema del Mercosur) contrasta con el modelo comunitario europeo, cuya lógica se desarrolla en la Unidad 5.

\section{Unidad 3 --- Especificidad del derecho de la integración}

La especificidad constituye la primera característica común a todo derecho de la integración: el órgano de un proceso de integración únicamente puede actuar en las materias y temas que le fueron específicamente asignados por el tratado constitutivo.

\begin{example}{La especificidad según el objeto del tratado}
Si el objeto de un proceso de integración es meramente económico (constituir un mercado común), el órgano solo puede regular el libre tránsito de los factores de producción; no podría, por ejemplo, dictar una norma sobre responsabilidad parental, porque esa materia no tiene relación con el mercado. La facultad del órgano comunitario se circunscribe a la competencia otorgada según el objeto del tratado.
\end{example}

\begin{quote}
\small
\textbf{Consulta:} ``Mi pregunta iba a esto de la especificidad, porque usted recién, al final, cuando hizo la pregunta, la mencionó relacionada con su tema.''

\textbf{Aclaración:} La especificidad es característica de cualquier derecho de la integración, con lo cual el derecho comunitario --- naturalmente --- también la tiene, pero no es una nota distintiva exclusiva suya: la comparte con cualquier proceso de integración, aunque este tenga menor nivel de compromiso institucional.
\end{quote}

\section{Unidad 4 --- Autonomía del derecho comunitario: los casos fundacionales del Tribunal}

Más allá de la especificidad, el derecho comunitario europeo desarrolló --- a través de la interpretación del Tribunal de Justicia de la Unión Europea --- características propias que no estaban escritas expresamente en los tratados originarios.

En una etapa muy temprana --- cuando ni siquiera existía todavía la actual Unión Europea, sino las Comunidades Europeas ---, el Tribunal se pronunció en dos sentencias fundacionales sobre la interpretación del derecho comunitario frente a una norma interna concreta de un Estado miembro (en uno de los casos, un tributo que un Estado pretendía cobrar en contradicción con el derecho comunitario).

\begin{example}{Los precedentes fundacionales del Tribunal}
En ambos casos, el Tribunal expuso principios que no estaban específicamente escritos en el tratado, pero que constituían la base misma del derecho comunitario. Estos precedentes corresponden a los conocidos casos fundacionales de la jurisprudencia europea sobre autonomía y primacía del derecho comunitario frente al derecho interno de los Estados miembros.
\end{example}

\begin{quote}
\itshape ``El Tribunal de la Unión Europea dice: goza también de autonomía [\ldots] el derecho comunitario no pertenece ni al derecho público, ni al derecho internacional, ni al derecho interno. Es derecho comunitario, es otra cosa.''
\end{quote}

% TODO: contenido incompleto o ambiguo en las notas originales --- el apunte señala que esta afirmación del Tribunal genera cierta reserva conceptual, sin desarrollar en qué consiste esa reserva.

Lo relevante es que el Tribunal fija la autonomía como principio: el derecho comunitario es, según su propia formulación, una ``tercera vía'' o ``tercer carril'', distinto tanto del derecho internacional como del derecho interno.

\begin{note}
La Comunidad Andina también cuenta con órganos supranacionales y replicó estas mismas características y principios del derecho comunitario europeo. Es, en el ámbito americano, lo más avanzado en materia de derecho de la integración: sus órganos son supranacionales, se expiden por mayoría y cuenta con un tribunal propio que ha seguido de cerca las pautas y parámetros desarrollados por la Unión Europea.
\end{note}

\section{Unidad 5 --- Supremacía y efecto directo (vertical y horizontal)}

A partir de la autonomía, el Tribunal fue estableciendo el resto de las características del derecho comunitario.

\subsection*{a) Supremacía}

La supremacía opera frente a las leyes de los Estados. Cuando una norma comunitaria entra en vigor, deja sin efecto automáticamente toda la normativa de derecho interno que la contradiga; el juez nacional debe inaplicar la norma interna contraria, porque esa norma está violando una norma superior. A la inversa, una vez dictada una norma comunitaria en determinado sentido, el Estado no puede dictar ninguna norma interna en sentido contrario.

\begin{quote}
\itshape ``Le marca el ámbito en el cual impide entrar en vigor de cualquier norma interna en sentido contrario y, al mismo tiempo, deroga de hecho y cesa la vigencia de las normas que están en sentido contrario.''
\end{quote}

\subsection*{b) Aplicabilidad inmediata}

Significa que el derecho comunitario no requiere ningún trámite previo de incorporación a la legislación de cada Estado: una vez terminado el procedimiento de sanción de la norma y su publicación oficial, la norma se aplica automáticamente en todos los Estados que participan de la Unión Europea, sin que sea necesario ningún procedimiento interno de incorporación.

\begin{note}
A diferencia del sistema de ``doble chequeo'' visto para el derecho derivado en otros procesos de integración (Unidad 2), la norma comunitaria no necesita ningún procedimiento de incorporación: entra en vigor para todos, incluso para el Estado que no la incorporó.
\end{note}

\subsection*{c) Efecto directo}

El efecto directo va de la mano con la aplicabilidad inmediata: la norma comunitaria no solo tiene vigencia inmediata, sino que además hace nacer --- en cabeza del ciudadano europeo --- un derecho o una obligación de forma automática, con independencia de que el Estado actúe o no para implementarla.

Si la norma exige que el Estado haga algo para que el ciudadano goce de un derecho y el Estado no lo hace, como el derecho ya está ``en su cabeza'', el ciudadano puede acudir directamente al Tribunal por inacción del Estado, mediante un recurso directo por incumplimiento.

El efecto directo se despliega en dos direcciones:
\begin{itemize}
    \item \textbf{Efecto directo vertical}: el ciudadano puede invocarlo frente al Estado.
    \item \textbf{Efecto directo horizontal}: el ciudadano puede invocarlo también frente a otros particulares, entre pares.
\end{itemize}

\begin{table}[H]
\centering
\caption{Características del derecho comunitario}
\begin{tabularx}{\textwidth}{@{} l X @{}}
\toprule
\textbf{Característica} & \textbf{Contenido central} \\
\midrule
Especificidad & Común a todo derecho de la integración: el órgano solo puede actuar en las materias que le fueron asignadas por el tratado. \\
Autonomía & No pertenece ni al derecho público, ni al internacional, ni al interno; es un ``tercer carril'' jurídico. \\
Supremacía & Deroga o deja sin efecto automáticamente cualquier norma interna contraria, en ambos sentidos temporales. \\
Aplicabilidad inmediata & No requiere procedimiento interno de incorporación; rige desde la publicación oficial. \\
Efecto directo (vertical y horizontal) & Crea derechos u obligaciones directamente en cabeza del ciudadano, oponibles al Estado y entre particulares. \\
\bottomrule
\end{tabularx}
\end{table}

\begin{quote}
\small
\textbf{Consulta:} ¿Cómo se concilian la autonomía --- que sostiene que el derecho comunitario es distinto y no pertenece a ningún otro orden --- con la supremacía, si en apariencia van por carriles distintos?

\textbf{Aclaración:} Es una tensión real, que se explica porque la Unión Europea no se comporta de la misma manera en todos los ámbitos ni con todo tipo de normas: según la materia, a veces rige la cooperación y no siempre la supralegalidad se aplica del mismo modo. Además, los distintos tipos de norma comunitaria --- desarrollados en el Bloque 2 --- tienen distinto grado de fuerza jurídica, lo que también incide en cómo operan estos principios en cada caso concreto.
\end{quote}

% =========================================================================
% BLOQUE 2
% =========================================================================
\chapter{Fuentes derivadas y armonización legislativa}

\section{Unidad 6 --- Tipos de normas del derecho derivado: reglamento, decisión y directiva}

Así como en el derecho interno argentino existen distintos tipos de normas según el órgano que las dicta --- una sentencia judicial, obligatoria solo para las partes de un caso concreto; una ley del Congreso, de alcance general ---, en la Unión Europea, sobre todo después del Tratado de Lisboa (cuando el Consejo colegisla con el Parlamento), los órganos pueden dictar distintos tipos de normas, cada una con un efecto jurídico específico.

\subsection*{1) El reglamento}

Es una norma general, que se aplica a todos los Estados y a todos los ciudadanos desde el momento en que entra en vigor. Goza plenamente de supremacía, efecto directo y aplicabilidad inmediata. Impide la entrada en vigor de cualquier norma nacional contraria y deroga o deja sin efecto cualquier norma interna en sentido contrario.

\begin{quote}
\itshape ``Fuerte el reglamento. Es una norma fuerte.''
\end{quote}

\subsection*{2) La decisión}

También tiene efecto directo, aplicabilidad inmediata y supremacía, pero --- a diferencia del reglamento --- es una norma particular: se dicta para resolver un problema puntual planteado por determinado grupo de Estados, y se aplica únicamente a las partes vinculadas en ese asunto concreto.

\begin{example}{Emergencia volcánica}
Ante una emergencia propia de ciertos Estados --- por ejemplo, una emergencia volcánica que no afecta a los demás ---, corresponde dictar una decisión y no un reglamento: ``no tengo ningún volcán, ¿por qué me vas a aplicar [la norma] si en mi territorio no existe este problema?''. Que la decisión sea ``particular'' no significa que se dicte para una persona en especial, sino para el grupo de países que comparten ese problema específico.
\end{example}

\subsection*{3) La directiva}

Fija el objetivo o las pautas generales que se quiere lograr, pero deja en manos del Estado la potestad de decidir cómo implementarlo.

\begin{quote}
\itshape ``Te obliga al fin, pero no al medio.''
\end{quote}

La directiva no tiene aplicabilidad inmediata, porque requiere que el Estado haga algo para ponerla en funcionamiento, pero sí tiene efecto directo: desde su entrada en vigor ya hace nacer en cabeza del ciudadano un derecho que lo habilita a reclamar frente al Estado si este no adecuó su normativa interna, mediante el recurso directo ante el Tribunal por incumplimiento del Estado.

\begin{example}{Igualdad de trato en el ingreso universitario}
Si una directiva establece ``igualdad de trato en el acceso a la educación'' y un Estado mantiene una norma interna que reserva el ingreso universitario solo a sus nacionales, sin derogarla, el ciudadano de otro Estado miembro puede invocar el recurso directo para que se obligue al Estado incumplidor a poner en funcionamiento esa igualdad de trato. ``Igualdad de trato'' no implica necesariamente ``educación gratuita para todos'': cada Estado puede optar por el modo concreto de garantizar esa igualdad, según sus posibilidades presupuestarias, su idiosincrasia, sus costumbres o su religión.
\end{example}

Esta figura más flexible existe porque el órgano o la autoridad que debe dar solución a un tema es la que está más próxima a ese tema: ciertas materias conviene dejarlas a criterio de cada Estado, para que la ciudadanía no sienta que la Unión Europea la absorbió por completo.

\begin{table}[H]
\centering
\caption{Reglamento, decisión y directiva}
\begin{tabularx}{\textwidth}{@{} l X X X X @{}}
\toprule
\textbf{Norma} & \textbf{Alcance} & \textbf{Aplicabilidad inmediata} & \textbf{Efecto directo} & \textbf{Contenido} \\
\midrule
Reglamento & General (todos los Estados y ciudadanos) & Sí & Sí & Regula la solución de fondo \\
Decisión & Particular (Estados o partes involucradas) & Sí & Sí & Resuelve un caso o problema puntual \\
Directiva & General en el objetivo & No (requiere acto del Estado) & Sí & Fija el fin, deja libre el medio \\
\bottomrule
\end{tabularx}
\end{table}

\section{Unidad 7 --- Procedimientos de armonización legislativa: uniformidad, armonía y aproximación}

No existe una única forma de armonizar las legislaciones de los Estados que integran un proceso de integración; pueden distinguirse tres procedimientos.

\subsection*{1) Uniformidad legislativa}

Es el procedimiento más sencillo: se dicta una norma que impone la misma solución de fondo para todos los Estados.

\begin{example}{La analogía del uniforme escolar}
Del mismo modo en que un uniforme escolar hace que todos los alumnos sean iguales, si se regula por uniformidad en materia laboral se establece que todos los trabajadores --- franceses, españoles, etc. --- tendrán la jornada laboral limitada a 8 horas y vacaciones mínimas de 15 días.
\end{example}

La solución de fondo es una sola, y por eso, cualquiera sea el juez que intervenga --- portugués, español o francés ---, el resultado del caso será el mismo, porque se aplica una única norma con una única interpretación.

\subsection*{2) Armonía legislativa}

En este procedimiento no se impone una única solución de fondo, sino que se elige un criterio común --- por ejemplo, un mismo punto de conexión --- para determinar qué ley se va a aplicar en cada caso, dejando que la solución sustancial siga siendo, eventualmente, distinta según el Estado.

\begin{example}{Ingreso universitario}
Puede acordarse como punto de conexión el domicilio del solicitante (criterio común para todos los Estados), pero cada Estado conserva la potestad de fijar sus propios requisitos de ingreso: unos toman examen, otros no.
\end{example}

\begin{example}{Contratos internacionales}
La armonía podría fijar como criterio uniforme que se aplique siempre ``la ley del lugar de cumplimiento del contrato''. Así, todos los Estados aplican la misma regla de asignación de competencia legislativa, aunque la solución de fondo --- los derechos y obligaciones del comprador y del vendedor --- termine siendo distinta según el derecho de fondo de cada país donde efectivamente se cumpla el contrato.
\end{example}

\subsection*{3) Aproximación legislativa}

Es un procedimiento de armonización más espontáneo que los dos anteriores: no responde a una decisión política intencional de los Estados reunidos con el objeto expreso de armonizar, sino que se produce naturalmente, como consecuencia de la propia evolución social y cultural de cada Estado.

\begin{example}{La progresiva admisión del divorcio vincular}
La progresiva admisión del divorcio vincular en distintos países se fue extendiendo Estado por Estado a medida que la sociedad y la cultura evolucionaban, sin que mediara un acuerdo previo entre los Estados para lograrlo. La armonización, en este caso, fue más una evolución de la sociedad y de la cultura: se dio espontáneamente, no obligatoriamente.
\end{example}

\begin{note}
La Unión Europea utiliza formalmente la denominación ``aproximación legislativa'' para su propio método de armonización, aunque en rigor conceptual eso no es exactamente lo que hace: recurre a la uniformidad o a la armonía según la materia de que se trate, pero prefirió un nombre más suave para no dar la sensación de imponer soluciones por la fuerza.
\end{note}

\section{Unidad 8 --- Caso aplicado: el Reglamento europeo sobre insolvencia}

El Reglamento de la Unión Europea sobre insolvencia (procedimientos de quiebra) demuestra que la elección entre reglamento, decisión y directiva (tipo de norma) y uniformidad, armonía y aproximación (método de armonización) constituyen dos ejes independientes que pueden combinarse de distintas maneras.

Se trata de un reglamento --- por lo tanto, formalmente goza de supremacía, efecto directo y aplicabilidad inmediata ---, pero que en su contenido de fondo aplica el método de la armonía y no el de la uniformidad. El reglamento remite la ley aplicable al lugar donde se encuentra el centro de intereses principales del deudor: según dónde esté ese centro, la ley de quiebra concretamente aplicable será distinta.

Lo que el reglamento sí unifica --- por vía de armonía --- son ciertos principios y mecanismos generales que todos los Estados deben respetar:

\begin{itemize}
    \item El deber de colaboración entre la quiebra principal y los procedimientos secundarios: quien quiebra en un Estado queda inhabilitado para funcionar en los demás.
    \item La posibilidad de que el síndico del procedimiento principal intervenga en los procedimientos secundarios en representación de los acreedores del proceso principal.
    \item La prohibición de discriminar entre acreedores nacionales y extranjeros.
\end{itemize}

\begin{example}{Contraste con el derecho argentino}
La ley argentina de quiebras dispone que los acreedores extranjeros cobran sobre el saldo, lo cual sería incompatible con el principio de igualdad de trato entre acreedores que exige el reglamento europeo: si no puede discriminarse entre acreedores nacionales y extranjeros, no podría aplicarse una normativa donde los extranjeros cobran sobre el saldo, como ocurre en la ley de quiebras argentina.
\end{example}

\begin{quote}
\itshape ``Era imposible que nos pongamos de acuerdo en un tema tan complicado como insolvencia'', dado que, al estar integradas las economías de los Estados, la baja de un centro productivo en un lugar genera una cadena de efectos e incumplimientos en otros.
\end{quote}

Ante la imposibilidad de unificar la solución de fondo, se optó por establecer, mediante reglamento, unas reglas mínimas comunes de cooperación y no discriminación, dejando el resto librado al derecho interno de cada Estado según el criterio de conexión fijado.

\section{Unidad 9 --- Síntesis: relación de género a especie con el derecho de la integración}

El derecho comunitario y el derecho de la integración en general se relacionan como género y especie: el derecho comunitario es una especie dentro del género más amplio ``derecho de la integración''.

Por ello, comparte con todo derecho de la integración ciertos principios comunes --- la especificidad, la igualdad de trato ---, pero además presenta rasgos propios y más intensos --- autonomía, supremacía, aplicabilidad inmediata, efecto directo vertical y horizontal --- que no necesariamente están presentes en otros procesos de integración de menor grado de compromiso institucional.

% =========================================================================
% BLOQUE 3
% =========================================================================
\chapter{Instituciones, tensiones actuales y Mercosur en perspectiva comparada}
\markboth{Bloque 3. Instituciones y Mercosur}{Bloque 3. Instituciones y Mercosur}

\section{Unidad 10 --- El Parlamento Europeo: composición y sistema de votación}

El Parlamento del Mercosur mantiene funciones eminentemente consultivas y de propuesta, sin mayor variación respecto de las que tenía anteriormente la Comisión Parlamentaria Conjunta, y no ejerce funciones de codecisión legislativa.

Esto contrasta fuertemente con el Parlamento Europeo, integrado por diputados elegidos en representación directa de los ciudadanos de cada Estado miembro. La distribución de bancas responde a un sistema de proporcionalidad decreciente (denominada también ``proporcionalidad inversa''): a mayor población de un Estado, mayor cantidad de escaños le corresponden, pero cada diputado de un Estado más poblado termina representando a más ciudadanos que un diputado de un Estado menos poblado.

\begin{example}{Límite físico de bancas}
Existe un límite físico de bancas en la sala (750), por lo que ningún Estado puede superar un determinado tope de representantes (aproximadamente 95), aun cuando --- de aplicarse una proporcionalidad estrictamente directa --- le hubiera correspondido un número mayor (a modo de ilustración, se menciona la cifra de 120).

\textit{``¿Por qué digo inversa? Porque esa es directa, que la cantidad de representantes que haya represente a tal cantidad de población: por cada uno van a representar a 10.000 ciudadanos. Digo inversa porque, mientras más población haya, ese diputado representa a más cantidad de ciudadanos.''}

Así, un diputado de un país muy poblado termina representando, por ejemplo, a 15.000 ciudadanos, mientras que uno de un país poco poblado representa a solo 5.000. Por eso se llama ``inversa'' o ``decreciente'': el peso relativo de cada banca decrece a medida que aumenta la población del Estado que representa.
\end{example}

\subsection*{La doble mayoría}

Sobre el sistema de votación, existe una evolución desde la mayoría simple hacia la mayoría doble (calificada) para determinados temas. La doble mayoría exige, en forma acumulativa:

\begin{itemize}
    \item El 55\% de los Estados que voten a favor, y
    \item Que esos Estados representen, en conjunto, al 65\% de la población total de la Unión.
\end{itemize}

\begin{example}{Doble mayoría no alcanzada}
Si se alcanza el 55\% de los Estados a favor, pero esos Estados son mayoritariamente de escasa población --- y los que se oponen son los más poblados ---, puede no alcanzarse el 65\% de población requerido (por ejemplo, llegar solo al 62\%). En ese caso, la propuesta no se aprueba, porque no se cumple el requisito de doble mayoría exigido para ese tipo de norma.
\end{example}

Cuando se pregunta por qué otros procesos de integración ``no son como la Unión Europea'', conviene recordar que a la Unión Europea ``le faltan 2000 años de historia'' --- es decir, un largo recorrido evolutivo e institucional ---, y que, más allá de ese proceso histórico, lo decisivo fue una firme decisión política de sostener el camino de la integración de manera irreversible.

\section{Unidad 11 --- Soberanía, globalización y el debate migratorio actual en Europa}

Existe una tensión entre ceder soberanía en favor de la integración y preservar la individualidad como Estado soberano. La especificidad (Unidad 3) resulta fundamental precisamente porque permite al Estado seguir conservando su individualidad soberana en todo aquello que no delegó expresamente.

Si un Estado sigue cediendo competencias de manera indefinida, se profundiza el fenómeno de la globalización, lo que ha derivado en un nuevo reclamo ciudadano en Europa: la sensación de tener menos derechos frente a quienes llegan de otros Estados, en un contexto de fuerte debate migratorio.

\begin{example}{La casa y el invitado}
A modo de ilustración pedagógica, se plantea la situación de abrir la propia casa a un invitado. El anfitrión puede recibirlo con los brazos abiertos, pero no puede pretender que el invitado le imponga sus propias costumbres dentro de su propia casa (por ejemplo, en materia de mascotas o de fumar dentro de la vivienda).

\textit{``No puedo pretender que el otro adecúe su casa a mis necesidades. Si no, no voy: voy a otro que responda lo que yo quiero.''}

Trasladada al plano estatal, la analogía sugiere que quien decide instalarse en otro país debe acomodarse a la idiosincrasia de ese Estado, sin que ello implique dejar de respetar la propia cultura de origen, pero tampoco pretender imponerla al Estado receptor: ``Yo puedo respetar la tuya, pero acomodate vos a la nuestra. No es que yo tengo que dejar de actuar como era. Sos vos el que viniste.''
\end{example}

Este tipo de tensiones --- que parecen obvias en el plano de la vida doméstica --- no fueron suficientemente valoradas por los Estados en el plano de la integración regional, y hoy, frente a fenómenos de migración masiva, los Estados europeos se preguntan cómo ``frenar'' procesos que en su momento impulsaron sin advertir todas sus consecuencias futuras.

\begin{example}{Matrimonio igualitario y adopción en España}
Como ilustración de una discusión análoga entre evolución social y consecuencia lógica no prevista, se menciona el debate que se dio en España a partir de la sanción del matrimonio igualitario y las consecuencias posteriores en materia de adopción, que constituían la consecuencia lógica de la propia decisión legislativa original, aunque en su momento no fueron dimensionadas por completo por la sociedad.
\end{example}

\begin{note}
Este tipo de temas se introducen deliberadamente como recurso didáctico para estimular el pensamiento crítico sobre las consecuencias de largo plazo de las decisiones de integración.
\end{note}

\section{Unidad 12 --- El Mercosur como unión aduanera imperfecta o incompleta}

\begin{quote}
\small
\textbf{Consulta:} En el marco de la ALADI se vieron muchos tratados: ¿todos fueron de libre comercio?

\textbf{Aclaración:} No todos: hay acuerdos de alcance parcial y de libre comercio, y hay acuerdos de integración. El Mercosur es un acuerdo de integración, al igual que la Comunidad Andina, y ambos se encuentran, a su vez, dentro del marco institucional de la ALADI.
\end{quote}

El Mercosur se clasifica como una unión aduanera ``imperfecta'' o, según una terminología preferible, ``incompleta'':

\begin{quote}
\itshape ``A mí no me gusta decir imperfecto, porque decir que algo es imperfecto es decir que está mal [\ldots] no estar completo significa que hiciste algunos acuerdos de aranceles en algunos espacios, pero todavía no llegaste al arancel externo común en la totalidad de bienes y servicios que circulan en el área.''
\end{quote}

El Mercosur avanzó en algunas rebajas arancelarias intrazona, pero aún no alcanzó un arancel externo común pleno para la totalidad de bienes y servicios, y persisten fenómenos como el contrabando en los rubros donde la barrera arancelaria todavía no fue levantada.

Asimismo, dentro del objetivo de mercado común --- libre tránsito de mercaderías, bienes, servicios y capitales ---, el Mercosur presenta una regulación parcial o inexistente de algunos factores de producción:

\begin{itemize}
    \item Hay avances parciales en materia de libre circulación de personas (por ejemplo, patente única y documento único).
    \item Existe cierta normativa en materia laboral.
    \item No existe todavía una ley aplicable uniforme en materia de contratos dentro del ámbito del Mercosur (compraventa de mercaderías, prestación de servicios), a pesar de tratarse de un factor de producción relevante.
\end{itemize}

\begin{important}
El Mercosur es, entonces, una zona de libre comercio intrazona incompleta, en camino hacia un mercado común que todavía no se ha completado en su totalidad. Es preferible el término ``incompleto'' por sobre ``imperfecto'', precisamente para no cargar de un juicio de valor negativo a un proceso que se encuentra en desarrollo.
\end{important}

% =========================================================================
% CORNELL
% =========================================================================
\clearpage
\chapter*{Resumen Cornell}
\addcontentsline{toc}{chapter}{Resumen Cornell}
\markboth{Resumen Cornell}{Resumen Cornell}

La siguiente tabla organiza, en formato Cornell, las preguntas o conceptos clave de cada unidad junto con su desarrollo, a modo de repaso rápido antes de la evaluación.

\begin{longtable}{
    >{\RaggedRight\arraybackslash}p{0.32\textwidth}
    >{\RaggedRight\arraybackslash}p{0.58\textwidth}
}
\toprule
\textbf{Preguntas / conceptos clave} & \textbf{Notas / respuesta / explicación} \\
\midrule
\endfirsthead

\toprule
\textbf{Preguntas / conceptos clave} & \textbf{Notas / respuesta / explicación} \\
\midrule
\endhead

\textbf{¿Qué jerarquía tienen los tratados de integración frente a la ley interna en la Argentina?} &
Desde el caso ``Cafés La Virginia'' y luego con la reforma de 1994 (art. 75, inc. 24, CN), los tratados de integración tienen jerarquía superior a las leyes. Esto es una decisión de derecho interno argentino, no una regla universal del derecho de la integración. \\
\addlinespace
\cmidrule{1-2}

\textbf{¿Cómo entra en vigor una norma de derecho derivado en sistemas con ``doble chequeo'' (por ejemplo, el Mercosur)?} &
La norma se notifica a través de la secretaría; cada Estado la incorpora según su propio sistema interno y comunica esa incorporación; recién cuando todos los Estados incorporaron, la secretaría informa la entrada en vigencia para todos. \\
\addlinespace
\cmidrule{1-2}

\textbf{¿Qué significa la especificidad del derecho de la integración?} &
El órgano de un proceso de integración solo puede actuar en las materias que le fueron específicamente asignadas por el tratado constitutivo. Es una característica común a todo derecho de la integración, no exclusiva del derecho comunitario. \\
\addlinespace
\cmidrule{1-2}

\textbf{¿Por qué el derecho comunitario es autónomo?} &
El Tribunal de Justicia de la Unión Europea estableció, en sus casos fundacionales, que el derecho comunitario no pertenece ni al derecho público, ni al internacional, ni al interno: es un ``tercer carril'' jurídico propio. \\
\addlinespace
\cmidrule{1-2}

\textbf{¿Qué produce la supremacía del derecho comunitario?} &
Deja sin efecto automáticamente cualquier norma interna contraria (deroga hacia atrás e impide dictar normas contrarias hacia adelante). El juez nacional debe inaplicar la norma interna que contradiga la norma comunitaria. \\
\addlinespace
\cmidrule{1-2}

\textbf{¿Qué diferencia hay entre aplicabilidad inmediata y efecto directo?} &
Aplicabilidad inmediata: la norma rige automáticamente en todos los Estados desde su publicación, sin trámite de incorporación. Efecto directo: la norma crea, además, derechos u obligaciones en cabeza del ciudadano, invocables frente al Estado (vertical) o frente a otros particulares (horizontal). \\
\addlinespace
\cmidrule{1-2}

\textbf{¿Qué diferencia hay entre reglamento, decisión y directiva?} &
Reglamento: norma general para todos los Estados y ciudadanos. Decisión: norma particular para el grupo de Estados afectados por un problema puntual. Directiva: fija el fin obligatorio, mas deja libre el medio de implementación al Estado. \\
\addlinespace
\cmidrule{1-2}

\textbf{¿Cómo se armonizan las legislaciones nacionales?} &
Por uniformidad (misma solución de fondo para todos), por armonía (mismo criterio o punto de conexión, con solución de fondo eventualmente distinta) o por aproximación (evolución espontánea, no decidida políticamente). \\
\addlinespace
\cmidrule{1-2}

\textbf{¿Por qué el Reglamento europeo de insolvencia aplica armonía y no uniformidad?} &
Porque unificar la solución de fondo en materia de insolvencia resultaba inviable dado el grado de integración económica; se optó por fijar reglas comunes de cooperación y no discriminación, remitiendo la ley de fondo al centro de intereses principales del deudor. \\
\addlinespace
\cmidrule{1-2}

\textbf{¿Cómo se relaciona el derecho comunitario con el derecho de la integración en general?} &
Es una relación de género a especie: el derecho comunitario es una especie dentro del género más amplio ``derecho de la integración'', y presenta rasgos propios (autonomía, supremacía, aplicabilidad inmediata, efecto directo) además de los comunes a todo proceso de integración. \\
\addlinespace
\cmidrule{1-2}

\textbf{¿Cómo se distribuyen los escaños y los votos en el Parlamento Europeo?} &
Por proporcionalidad decreciente (más población, más bancas, pero cada banca representa a más ciudadanos cuanto más poblado es el Estado, debido al límite de 750 bancas). Para ciertos temas se exige doble mayoría: 55\% de los Estados que representen, además, el 65\% de la población. \\
\addlinespace
\cmidrule{1-2}

\textbf{¿Qué tensión existe entre integración y soberanía?} &
Ceder competencias de manera indefinida profundiza la globalización y puede generar reclamos ciudadanos (por ejemplo, en el debate migratorio europeo). La especificidad protege la individualidad soberana del Estado en todo lo no delegado expresamente. \\
\addlinespace
\cmidrule{1-2}

\textbf{¿Por qué el Mercosur se clasifica como unión aduanera imperfecta o incompleta?} &
Porque todavía no alcanzó un arancel externo común pleno para la totalidad de bienes y servicios, y porque la libre circulación de los factores de producción (personas, servicios, contratos) está solo parcialmente regulada. \\
\addlinespace
\midrule

\multicolumn{2}{p{0.94\textwidth}}{
\textbf{Resumen:} La clase reconstruye el edificio conceptual del derecho de la integración a partir de una advertencia metodológica clave: que un Estado --- como la Argentina, desde el art. 75, inc. 24, de la Constitución Nacional --- reconozca jerarquía superior a los tratados de integración es una decisión de su propio derecho interno, y no una nota inherente a todo derecho de la integración. Sobre esa base, el derecho comunitario europeo se presenta como una especie particularmente evolucionada dentro del género ``derecho de la integración'': comparte con cualquier proceso de integración la especificidad, pero añade, por vía de la interpretación del Tribunal de Justicia, la autonomía, la supremacía, la aplicabilidad inmediata y el efecto directo (vertical y horizontal). Estas características se plasman de modo distinto según el tipo de norma --- reglamento, decisión o directiva --- y según el método de armonización elegido --- uniformidad, armonía o aproximación ---, como demuestra el caso del Reglamento europeo de insolvencia, fuerte en su forma pero flexible en su contenido de fondo. La clase sitúa además estos conceptos en un plano institucional y político más amplio: el diseño del Parlamento Europeo (proporcionalidad decreciente y doble mayoría) expresa un largo proceso histórico y una decisión política sostenida de integrarse, mientras que la tensión actual entre soberanía, globalización y migración recuerda que ceder competencias sin límite tiene un costo que las sociedades hoy discuten. Finalmente, la comparación con el Mercosur --- una unión aduanera todavía incompleta, sin arancel externo común pleno ni libre circulación total de los factores de producción --- permite dimensionar cuánto camino institucional resta recorrer en los procesos de integración regional latinoamericanos respecto del modelo europeo.
} \\
\bottomrule
\end{longtable}

\end{document}
